\clearpage
\chapter*{第190次报告勘误}
\addcontentsline{toc}{chapter}{第190次报告勘误}
\sid{e190}
\setcounter{footnote}{0}
\src{300}

\noindent\textbf{3. 下降技术与代数几何中的存在性定理 [TDTE I]。
I：一般理论。沿忠实平坦态射的下降。}

\begin{center}
[布尔巴基讨论班，第12卷，1959/60年度，第190号，29页。]
\end{center}

\noindent\textbf{第1页，第10行。} --- 现在看来，说下降技术“是代数几何中
大多数存在性定理的基础”似乎过甚。对于 TDTE I 至 VI 系列前两篇报告所处理
的非射影技术，这在很大程度上成立；但对于射影技术（TDTE IV、V、VI）则不
成立。

\noindent\textbf{第5页，第16行。} --- 无需假设 $\alpha$ 是一个
$\mathcal F$-下降态射。

\noindent\textbf{第20页，评注。} --- 拟有限、平展且满射的态射
$S'\longrightarrow S$，或者态射
$\operatorname{Spec}(\widehat A)\longrightarrow\operatorname{Spec}(A)$，
未必是严格下降态射；即使 $A$ 是代数闭域 $k$ 上一条代数曲线的局部环，且
$S=\operatorname{Spec}(A)$，情形仍然如此。确实可以找到固有光滑态射
$f:X\longrightarrow S$，使 $X$ 成为一个以 $S$ 为基、以椭圆曲线 $E$ 为
结构群的主丛，并且 $f':X'\longrightarrow S'$ 射影而 $f$ 不射影。因此，这同时
也是阿贝尔概形作用下非等平凡主齐性丛的一个例子。

\noindent\textbf{第26页，第9行。} --- 应读作“Chow--Lang”，而非
“Igusa--Lang”。

关于下降理论的各项细节，见 SGA VI、VII、VIII。
